\subsection{Subspace tracking}
One problem that arises when working with principal component analysis is that the result, which is an covariance matrix, is often of higher dimension and computationally challenging to diagonalize directly via a full eignedecomposition. The subspace iteration method solves this issue by iteratively estimating the dominant eigenvectors instead of solving the eigenproblem at each iteration step. \textcite[p.~270]{Oja1982SimplifiedNM} proves that a simple Hebbian-type update rule combined with normalization converges with probability one to the dominant eigenvector of the input correlation matrix, provided the largest eigenvalue has multiplicity one. This provides the theoretical basis for implementing the subspace iteration method to solve the eigenvector problem.
Following this principle the eigenvector solver used in the principal component analysis routine follows the classical power iteration scheme described by \textcite[p.~2932]{1468483}, which alternates a data compression step—multiplying the correlation matrix by the current subspace estimate—with an orthonormalization step at each iteration.
However \textcite[p.~2932]{1468483} points out that the classical power iteration method is unsuitable for real-time processing due to its computational cost, which scales with the ambient dimension $n$ of the data; this motivated the development of faster approximations such as API and FAPI. 
To achieve fast convergence requires the covariance matrix to be kept as small as possible. Since in this thesis the covariance matrix is computed over the $d$-dimensional embedding space rather than on an large continuous array, the power iteration remains computationally stable for real time use without requiring these approximations.

